\subsection{Criterios para que los conjuntos
$U_{S_n}(\mathcal F)$ o $U_{C_n}(\mathcal F)$ sean abiertos.}
\label{IV.2.s6.11-fr}

\begin{lemma}[6.11.1]
\label{IV.2.s6.11.1-fr}
Sean $X$ un preesquema localmente noetheriano, $\mathcal F$ un
$\mathcal O_X$-Módulo coherente; entonces la función
$x\leadsto\operatorname{dim.proj}(\mathcal F_x)$ es semicontinua
superiormente en $X$.
\end{lemma}

\begin{proof}
Es posible limitarse al caso en que $X=\operatorname{Spec}(A)$ es el espectro de un
anillo noetheriano y $\mathcal F=\widetilde M$, donde $M$ es un $A$-módulo
finitamente generado. Supongamos que, para un $x\in X$, se tiene
$\operatorname{dim.proj}(M_x)=n<+\infty$ (si $n=+\infty$, no hay nada que
demostrar); existe una resolución de $M$
\[
L_{n-1}\longrightarrow L_{n-2}\longrightarrow\ldots
\longrightarrow L_0\longrightarrow M\longrightarrow0
\]
\oldpage[IV-2]{159}
donde los $L_i$ son $A$-módulos libres \emph{de tipo finito} (pues $A$ es
noetheriano), de donde se obtiene una sucesión exacta
\[
0\longrightarrow R\longrightarrow L_{n-1}\longrightarrow L_{n-2}
\longrightarrow\ldots\longrightarrow L_0\longrightarrow M
\longrightarrow0
\]
donde $R$ es un $A$-módulo finitamente generado; se deduce una sucesión exacta
\[
0\longrightarrow R_x\longrightarrow(L_{n-1})_x\longrightarrow\ldots
\longrightarrow(L_0)_x\longrightarrow M_x\longrightarrow0
\]
donde los $(L_i)_x$ son $A_x$-módulos libres finitamente generados; como, por
hipótesis, $\operatorname{dim.proj}(M_x)=n$, esto implica que $R_x$ es un
$A_x$-módulo proyectivo finitamente generado (M, VI, 2.1) y, por tanto, un
$A_x$-módulo \emph{libre} de tipo finito (\textbf{0}$_{\rm III}$, 10.1.3). Se concluye que
existe una vecindad abierta $U$ de $x$ en $X$ tal que, para todo
$x'\in U$, $R_{x'}$ es un $A_{x'}$-módulo libre (\textbf{0}$_{\rm I}$, 5.2.7);
por consiguiente, $M_{x'}$ admite una resolución proyectiva de longitud $n$, es decir,
$\operatorname{dim.proj}(M_{x'})\leq n$, lo que prueba el lema.
\end{proof}
